\subsection{Базовые классы, их функции и использование}

Scikit-learn (sklearn) построен на основе единого, консистентного программного интерфейса приложений, основными
абстракциями которого являются \textbf{оценщики} (estimators), \textbf{преобразователи} (transformers) и
\textbf{предсказатели} (predictors).

\underline{Базовый класс \texttt{sklearn.base.BaseEstimator}}.
Является родительским классом для всех оценщиков.
Определяет общие методы:
\begin{itemize}
    \item \texttt{\_\_init\_\_()}: Конструктор для установки гиперпараметров
    \item \texttt{get\_params()}, \texttt{set\_params()}: Методы для получения и установки параметров модели
    (важно для настройки гиперпараметров и кросс-валидации)
\end{itemize}

\underline{Классы \texttt{sklearn.base.ClassifierMixin / RegressorMixin / TransformerMixin}}
Миксины (примеси), которые добавляют специфичное поведение:
\begin{itemize}
    \item \texttt{ClassifierMixin}: добавляет метод \texttt{score(X, y)}, возвращающий
    \textbf{accuracy} (долю правильных ответов) для классификаторов.
    \item \texttt{RegressorMixin}: добавляет метод \texttt{score(X, y)}, возвращающий
    \textbf{коэффициент детерминации $R^2$} для регрессоров.
    \item \texttt{TransformerMixin}: добавляет метод \texttt{fit\_transform(X[, y])}, который выполняет
    \texttt{fit()} и \texttt{transform()} за один шаг.
    Ускоряет написание кода.
\end{itemize}

\underline{Ключевые интерфейсы и методы}

Любой ``обученный'' объект в sklearn (модель, преобразователь) следует одному из двух основных шаблонов.

\textbf{Оценщик (Модель) для обучения с учителем} (\texttt{Classifier}, \texttt{Regressor})
\begin{itemize}
    \item \texttt{.fit(X\_train, y\_train)}.
    Основной метод для обучения модели на тренировочных данных.
    Изменяет внутреннее состояние объекта (становится ``обученным'').
    \item \texttt{.predict(X\_new)} использует обученную модель для предсказания целевой переменной на новых
    данных.
    Возвращает массив \texttt{y\_pred}.
    \item \texttt{.predict\_proba(X\_new)}: работает только для классификаторов и возвращает вероятности
    принадлежности к каждому классу.
    \item \texttt{.score(X, y)}: Возвращает метрику качества предсказаний на предоставленных данных
\end{itemize}

\textbf{Пример использования:}
\begin{lstlisting}[language=Python, label={lst:estimator}]
    # from sklearn.ensemble import RandomForestClassifier
    # from sklearn.datasets import load_iris
    # from sklearn.model_selection import train_test_split
    #
    # # Загрузка данных
    # X, y = load_iris(return_X_y=True)
    # X_train, X_test, y_train, y_test = train_test_split(X, y, test_size=0.3)
    #
    # # Создание и обучение оценщика
    # clf = RandomForestClassifier(n_estimators=100, random_state=42)
    # clf.fit(X_train, y_train) # Обучение (fit)
    #
    # # Использование обученной модели
    # y_pred = clf.predict(X_test) # Предсказание (predict)
    # accuracy = clf.score(X_test, y_test) # Оценка (score)
    # print(f"Accuracy: {accuracy:.2f}")
\end{lstlisting}

\textbf{Важное правило:} параметры преобразования (например, \texttt{scaler.mean\_}) всегда вычисляются
\textbf{только на тренировочной выборке} (\texttt{fit}), а затем применяются и к тренировочной, и к тестовой
(\texttt{transform}), чтобы избежать неучёта данных.

\subsection{Важность пайплайнов, пример}

\underline{Проблемы, которые решает Pipeline}

Без использования \texttt{Pipeline} процесс предобработки и обучения часто выглядит громоздким и ошибочным
\begin{enumerate}
    \item \textbf{Утечка данных (Data Leakage):} при ручном выполнении \texttt{fit\_transform} на всем
    датасете до разделения или при вычислении статистик (например, среднего для заполнения пропусков) на тесте.
    \item \textbf{Сложность воспроизведения:} много промежуточных шагов и переменных.
    \item \textbf{Неудобство кросс-валидации и поиска гиперпараметров:} невозможно корректно настроить
    гиперпараметры для этапа предобработки (например, степень полинома в \texttt{PolynomialFeatures}) совместно
    с параметрами модели в рамках GridSearchCV\@.
    \item \textbf{Сложность развёртывания:} пайплайн инкапсулирует всю цепочку преобразований и модель в
    один объект, который легко сохранить (\texttt{pickle} или \texttt{joblib}) и использовать в production.
    Без этого предстоит всё делать самому
\end{enumerate}

\underline{Пример создания и использования Pipeline}

\begin{lstlisting}[language=Python, label={lst:pipeline}, caption=Полный пример пайплайна с настройкой]
    # from sklearn.pipeline import Pipeline
    # from sklearn.compose import ColumnTransformer
    # from sklearn.preprocessing import StandardScaler, OneHotEncoder
    # from sklearn.impute import SimpleImputer
    # from sklearn.ensemble import RandomForestClassifier
    # from sklearn.model_selection import train_test_split, GridSearchCV
    # import pandas as pd
    # import numpy as np
    #
    # # 1. Создание синтетических данных
    # np.random.seed(42)
    # n_samples = 100
    # data = pd.DataFrame({
    #     'age': np.random.randint(18, 70, n_samples),
    #     'salary': np.random.normal(50000, 15000, n_samples),
    #     'city': np.random.choice(['Moscow', 'SPb', 'Kazan'], n_samples),
    #     'experience': np.random.normal(10, 4, n_samples),
    #     'target': np.random.choice([0, 1], n_samples)
    # })
    # # Добавим пропуски и выбросы для реалистичности
    # data.loc[10:15, 'salary'] = np.nan
    # data.loc[20, 'age'] = 150
    #
    # X = data.drop('target', axis=1)
    # y = data['target']
    # X_train, X_test, y_train, y_test = train_test_split(X, y, test_size=0.2, random_state=42)
    #
    # # 2. Определение этапов предобработки для разных типов признаков
    # numeric_features = ['age', 'salary', 'experience']
    # numeric_transformer = Pipeline(steps=[
    #     ('imputer', SimpleImputer(strategy='median')), # Заполняем пропуски медианой
    #     ('scaler', StandardScaler()) # Масштабируем
    # ])
    #
    # categorical_features = ['city']
    # categorical_transformer = Pipeline(steps=[
    #     ('imputer', SimpleImputer(strategy='constant', fill_value='missing')), # Заполняем пропуски константой
    #     ('onehot', OneHotEncoder(handle_unknown='ignore', sparse_output=False)) # One-Hot Encoding
    # ])
    #
    # # 3. Объединяем преобразователи в ColumnTransformer
    # preprocessor = ColumnTransformer(
    #     transformers=[
    #         ('num', numeric_transformer, numeric_features),
    #         ('cat', categorical_transformer, categorical_features)
    #     ])
    #
    # # 4. Создаем итоговый пайплайн: предобработка -> модель
    # full_pipeline = Pipeline(steps=[
    #     ('preprocessor', preprocessor),
    #     ('classifier', RandomForestClassifier(random_state=42))
    # ])
    #
    # # 5. Определяем сетку гиперпараметров для настройки
    # # Обратите внимание на синтаксис: имя этапа в пайплайне + __ + имя параметра
    # param_grid = {
    #     'preprocessor__num__imputer__strategy': ['mean', 'median'],
    #     'classifier__n_estimators': [50, 100, 200],
    #     'classifier__max_depth': [5, 10, None]
    # }
    #
    # # 6. Создаем и обучаем GridSearchCV с пайплайном
    # grid_search = GridSearchCV(full_pipeline, param_grid, cv=5, scoring='accuracy', n_jobs=-1)
    # grid_search.fit(X_train, y_train) # ВСЕ преобразования применяются корректно внутри каждого фолда CV!
    #
    # print(f"Лучшая точность: {grid_search.best_score_:.3f}")
    # print(f"Лучшие параметры: {grid_search.best_params_}")
    #
    # # 7. Оценка на тестовой выборке
    # best_model = grid_search.best_estimator_
    # test_score = best_model.score(X_test, y_test) # Пайплайн сам преобразует X_test!
    # print(f"Точность на тесте: {test_score:.3f}")
    #
    # # 8. Предсказание для новых данных
    # new_data = pd.DataFrame({
    #     'age': [25, 45],
    #     'salary': [60000, np.nan],
    #     'city': ['Moscow', 'Kazan'],
    #     'experience': [2, 15]
    # })
    # predictions = best_model.predict(new_data)
    # print(f"Предсказания: {predictions}")
\end{lstlisting}

\underline{Преимущества данного подхода}

\begin{itemize}
    \item \textbf{Защита от утечки:} \texttt{SimpleImputer} и \texttt{StandardScaler} вычисляют свои параметры
    \textbf{только} на тренировочном фолде во время кросс-валидации.
    \item \textbf{Чистота кода:} весь процесс упакован в один объект \texttt{best\_model}.
    \item \textbf{Удобная настройка:} можно подбирать гиперпараметры как для предобработки, так и для модели.
    \item \textbf{Простое развертывание:} объект \texttt{best\_model} можно сразу сохранить (\texttt{joblib.dump})
    и загрузить для предсказаний на новых сырых данных без необходимости помнить о порядке преобразований.
\end{itemize}